\subsection{History and relation to earlier work}
Reconstructing a sequence from corrupted copies has a broad information-theoretic
history; Levenshtein studied reconstruction under several combinatorial and
probabilistic channel models~\cite{Levenshtein2001}. Batu, Kannan, Khanna,
and McGregor formulated reconstruction from independent random deletion
traces in theoretical computer science, motivated in part by sequence
alignment~\cite{BKKM2004}. The missing positions make the problem different
from coordinatewise denoising: the same trace rank can originate at many
input positions.

For a fixed deletion probability, Holenstein, Mitzenmacher, Panigrahy, and
Wieder obtained a worst-case sample bound
$\exp(\widetilde O(\sqrt n))$~\cite{HMPW2008}. In independent papers
published at STOC 2017, De, O'Donnell, and Servedio, and Nazarov and Peres,
obtained $\exp(O(n^{1/3}))$ bounds and matching limitations for methods
using only trace means~\cite{DOS2019,NP2017}. Their complex-analytic approach
connects deletion statistics to bounded-coefficient polynomials; analytic
inequalities of Borwein and Erd\'elyi are part of the background to this
method~\cite{BE1997}. The mean-based lower bound is a restriction on that
method, not on arbitrary trace algorithms.

Chase improved the general sample upper bound to
$\exp(O(n^{1/5}\log^5n))$~\cite{Chase2021Upper}. On the lower-bound side,
Holden and Lyons proved an almost-$n^{5/4}$ bound, and Chase strengthened
it to $\Omega(n^{3/2}/\log^7n)$ for fixed deletion
probability~\cite{HoldenLyons2020,Chase2021Lower}.
Burudgunte, Valiant, and Wang subsequently proved the sample bound
\[
 \exp\bigl(p^{-7/3}(\log_2n)^{c_0}\bigr)
\]
for an absolute $c_0$, using nonlocal low-order statistics and a multiscale
propagation argument~\cite[Theorem~43]{BVW2026}. This bound is
quasipolynomial when $p^{-1}$ is polylogarithmic in $n$.
Theorem~\ref{thm:main}
builds on this strategy and provides an explicit uniform reconstruction
algorithm with the displayed bit-complexity bound.

Sample complexity and reconstruction time require different arguments.
At fixed confidence, maximum likelihood achieves worst-case reconstruction
using at most $O(n)$ times the optimal sample count, but evaluating it over
all candidate strings entails a large search~\cite[Corollary~1.1 of the
full version]{CGLSZ2024}. Burudgunte, Valiant, and Wang
explicitly distinguish testing two known strings from reconstructing an
unknown string, and leave quasipolynomial-time reconstruction open for
their approach~\cite[Section~1]{BVW2026}. The additional separation proved
here concerns every positive semidefinite moment completion of a wrong
prefix. It enables constructive feasibility tests in place of a search
over complete strings. Earlier polynomial-time results already cover
low-deletion regimes~\cite{CDLSS2021LowDeletion}.

The use of nonlocal statistics is essential to the comparison with known
barriers. Chen, De, Lee, and Servedio proved strong lower bounds for
statistical queries restricted to short \emph{consecutive} trace
blocks~\cite{CDLS2024}. Our statistics involve few ranks but may span the
whole trace, as in~\cite{BVW2026}. The generating-function identity used
to transfer them from input coordinates has antecedents in higher-order
trace reconstruction~\cite{CDLSS2021Smoothed,Chase2021Upper}.

The positive moment relaxation is a Boolean instance of the semidefinite
moment and sum-of-squares framework~\cite{Lasserre2001,Parrilo2000}.
The finite feasibility iteration is related to the classical relaxation
methods of Agmon and Motzkin--Schoenberg~\cite{Agmon1954,MotzkinSchoenberg1954}.
We prove the analytic separation and the rational implementation directly;
no external trace-reconstruction or optimization theorem is needed.
